% !Mode:: "TeX:UTF-8"
% ============================================================
% CONCLUSION
% ============================================================
\chapter*{Conclusion}
\addcontentsline{toc}{chapter}{Conclusion}
\markboth{CONCLUSION}{CONCLUSION}

This study guide has presented the theory, worked examples and exercises
needed for the practical part of a course in linear algebra. The material
starts with matrices and determinants and proceeds through exact methods
for systems of linear equations (Cramer’s rule, the matrix method, the
Gauss and Gauss–Jordan methods) and approximate ones (simple iteration and
the Gauss–Seidel method), vector algebra, elements of analytic geometry,
linear and Euclidean spaces, linear operators, eigenvalues and
eigenvectors, and quadratic forms.

Every topic has been connected with its use in information security: the
Hill cipher, error-correcting codes, secret sharing schemes, linear
cryptanalysis, lattice-based cryptography, network analysis and anomaly
detection. These applications are meant to show students that linear
algebra is not only an abstract theory but a tool they will use directly in
their profession. The short Python, Maple and Mathcad examples show how the
same computations are carried out in practice, where matrices are large and
hand calculation is no longer an option.

The more than two thousand exercises, all supplied with answers, give
students ample material for independent work and allow teachers to prepare
individual assignments. The knowledge and skills acquired here form the
basis for subsequent courses such as discrete mathematics, probability
theory, numerical methods, cryptography and coding theory.
